% GENERATED FILE. Edit canonical lessons, then run npm run generate:books.
% canonical-source-hash: fnv1a64:4158f2df7444d44d
% canonical-lessons: LA-C84-arca, LA-C84-charta, LA-C84-stilus, LA-C84-tabula, LA-C84-speculum

\chapter{Chest, Paper, Pen, Tablet, Mirror}
\label{ch:la-chest-paper-pen-tablet-mirror}
\hlchaptermodality{84}

\begin{chapteropening}
\textbf{By the end of this chapter:} \emph{I can name a chest, paper, a writing stick, a tablet and a mirror.}

Complete the last lesson of chapter 84: I can name a chest, paper, a writing stick, a tablet and a mirror.
\end{chapteropening}

\section[arca, arcae]{arca, arcae — a chest, a box}

\label{lesson:LA-C84-arca}



\begin{quote}
Before the new one: say the Latin for a cup, then the Latin for a knife.
\end{quote}

\subsection*{You'll want to know: arca, arcae}
\textbf{arca, arcae} — \textquotedblleft{}a chest, a box\textquotedblright{}.

\begin{cousinweb}
\textbf{arca} is usually linked with \emph{arceō}, \textquotedblleft{}to shut in, to keep off\textquotedblright{} (Lewis \& Short): a chest keeps things shut away. English \emph{ark} comes from \emph{arca}.
\end{cousinweb}

\begin{grammarlens}[title={What you've built}]
One more word for this chapter.
\end{grammarlens}

\subsection*{Your turn}
\begin{itemize}
\raggedright
  \item \emph{Say it:} \emph{arca}
  \item \emph{Say it:} \emph{arca}, once more
  \item \emph{Say it:} say \emph{arca}
  \item \emph{Recall:} say the Latin for a cup, then the Latin for a knife, then say \emph{arca} again
\end{itemize}

\begin{tcolorbox}[breakable,colback=teal!4,colframe=teal!35!black,title={Before you move on}]
What does arca, arcae mean? (\textquotedblleft{}A chest, a box\textquotedblright{}.) Say it once more.
\end{tcolorbox}
\section[charta, chartae]{charta, chartae — paper, a sheet}

\label{lesson:LA-C84-charta}



\begin{quote}
Before the new one: say the Latin for a knife, then the Latin for a chest.
\end{quote}

\subsection*{You'll want to know: charta, chartae}
\textbf{charta, chartae} — \textquotedblleft{}paper, a sheet\textquotedblright{}. English \emph{chart} and \emph{card}.

\begin{cousinweb}
\textbf{charta} was borrowed from Greek \emph{khartēs}, \textquotedblleft{}a sheet of papyrus\textquotedblright{} (Lewis \& Short), which may come from \emph{kharassō}, \textquotedblleft{}to scratch, to engrave\textquotedblright{} (EtymDB 2.1).
\end{cousinweb}

\begin{grammarlens}[title={What you've built}]
One more word for this chapter.
\end{grammarlens}

\subsection*{Your turn}
\begin{itemize}
\raggedright
  \item \emph{Say it:} \emph{charta}
  \item \emph{Say it:} \emph{charta}, once more
  \item \emph{Say it:} say \emph{charta}
  \item \emph{Recall:} say the Latin for a knife, then the Latin for a chest, then say \emph{charta} again
\end{itemize}

\begin{tcolorbox}[breakable,colback=teal!4,colframe=teal!35!black,title={Before you move on}]
What does charta, chartae mean? (\textquotedblleft{}Paper, a sheet\textquotedblright{}.) Say it once more.
\end{tcolorbox}
\section[stilus, stilī]{stilus, stilī — a stylus, a writing stick}

\label{lesson:LA-C84-stilus}



\begin{quote}
Before the new one: say the Latin for a chest, then the Latin for paper.
\end{quote}

\subsection*{You'll want to know: stilus, stilī}
\textbf{stilus, stilī} — \textquotedblleft{}a stylus, a writing stick\textquotedblright{}. English \emph{style}.

\begin{cousinweb}
\textbf{stilus} is probably inherited, from a Proto-Indo-European root \emph{*steyg-}, \textquotedblleft{}to prick, to pierce\textquotedblright{} (EtymDB 2.1); English \emph{stick} is a cousin (Lewis \& Short). The spelling with \emph{y} came later, from a mistaken link with Greek \emph{stylos}, \textquotedblleft{}pillar\textquotedblright{}.
\end{cousinweb}

\begin{grammarlens}[title={What you've built}]
One more word for this chapter.
\end{grammarlens}

\subsection*{Your turn}
\begin{itemize}
\raggedright
  \item \emph{Say it:} \emph{stilus}
  \item \emph{Say it:} \emph{stilus}, once more
  \item \emph{Say it:} say \emph{stilus}
  \item \emph{Recall:} say the Latin for a chest, then the Latin for paper, then say \emph{stilus} again
\end{itemize}

\begin{tcolorbox}[breakable,colback=teal!4,colframe=teal!35!black,title={Before you move on}]
What does stilus, stilī mean? (\textquotedblleft{}A stylus, a writing stick\textquotedblright{}.) Say it once more.
\end{tcolorbox}
\section[tabula, tabulae]{tabula, tabulae — a writing tablet, a board}

\label{lesson:LA-C84-tabula}



\begin{quote}
Before the new one: say the Latin for paper, then the Latin for a stylus.
\end{quote}

\subsection*{You'll want to know: tabula, tabulae}
\textbf{tabula, tabulae} — \textquotedblleft{}a writing tablet, a board\textquotedblright{}. English \emph{table}.

\begin{cousinweb}
\textbf{tabula} has no sure origin. Lewis \& Short connect it with \emph{taberna}, \textquotedblleft{}a booth, a shop\textquotedblright{}, but the link is uncertain.
\end{cousinweb}

\begin{grammarlens}[title={What you've built}]
One more word for this chapter.
\end{grammarlens}

\subsection*{Your turn}
\begin{itemize}
\raggedright
  \item \emph{Say it:} \emph{tabula}
  \item \emph{Say it:} \emph{tabula}, once more
  \item \emph{Say it:} say \emph{tabula}
  \item \emph{Recall:} say the Latin for paper, then the Latin for a stylus, then say \emph{tabula} again
\end{itemize}

\begin{tcolorbox}[breakable,colback=teal!4,colframe=teal!35!black,title={Before you move on}]
What does tabula, tabulae mean? (\textquotedblleft{}A writing tablet, a board\textquotedblright{}.) Say it once more.
\end{tcolorbox}
\section[speculum, speculī]{speculum, speculī — a mirror}

\label{lesson:LA-C84-speculum}



\begin{quote}
Before the new one: say the Latin for a stylus, then the Latin for a writing tablet.
\end{quote}

\subsection*{You'll want to know: speculum, speculī}
\textbf{speculum, speculī} — \textquotedblleft{}a mirror\textquotedblright{}.

\begin{cousinweb}
\textbf{speculum} is built on \emph{speciō}, \textquotedblleft{}to look at\textquotedblright{} (Lewis \& Short), with \emph{-ulum}, an ending for a tool: the thing one looks into. \emph{speciō} is usually traced to Proto-Indo-European \emph{*spek-}, \textquotedblleft{}to look\textquotedblright{} (EtymDB 2.1).
\end{cousinweb}

\begin{grammarlens}[title={What you've built}]
One more word for this chapter.
\end{grammarlens}

\subsection*{Your turn}
\begin{itemize}
\raggedright
  \item \emph{Say it:} \emph{speculum}
  \item \emph{Say it:} \emph{speculum}, once more
  \item \emph{Say it:} say \emph{speculum}
  \item \emph{Recall:} say the Latin for a stylus, then the Latin for a writing tablet, then say \emph{speculum} again
\end{itemize}

\begin{tcolorbox}[breakable,colback=teal!4,colframe=teal!35!black,title={Before you move on}]
What does speculum, speculī mean? (\textquotedblleft{}A mirror\textquotedblright{}.) Say it once more.
\end{tcolorbox}
